\documentclass[12pt,a4paper]{amsart}

\usepackage{amsmath,amssymb,amsthm}
\usepackage{mathtools}
\usepackage{hyperref}
\usepackage{geometry}
\geometry{margin=2.5cm}
\usepackage{booktabs}
\usepackage[ngerman]{babel}

% -------------------------------------------------------
% Theorem-Umgebungen
% -------------------------------------------------------
\newtheorem{theorem}{Satz}[section]
\newtheorem{lemma}[theorem]{Lemma}
\newtheorem{korollar}[theorem]{Korollar}
\newtheorem{proposition}[theorem]{Proposition}
\theoremstyle{definition}
\newtheorem{definition}[theorem]{Definition}
\newtheorem{bemerkung}[theorem]{Bemerkung}

% -------------------------------------------------------
% Eigene Befehle
% -------------------------------------------------------
\newcommand{\N}{\mathbb{N}}
\newcommand{\Z}{\mathbb{Z}}
\newcommand{\Q}{\mathbb{Q}}
\newcommand{\R}{\mathbb{R}}
\newcommand{\C}{\mathbb{C}}
\newcommand{\F}{\mathbb{F}}
\newcommand{\leg}[2]{\left(\frac{#1}{#2}\right)}
\DeclareMathOperator{\ord}{ord}
\DeclareMathOperator{\lcm}{kgV}
\DeclareMathOperator{\li}{li}
\DeclareMathOperator{\ggT}{ggT}

% -------------------------------------------------------
\begin{document}
% -------------------------------------------------------

\title{Brunnscher Satz: Konvergenz der Kehrwert-Reihe\\über Zwillingsprimzahlen}

\author{Michael Fuhrmann}
\address{Specialist Mathematics Project, Build 74}
\email{specialist-maths@localhost}
\date{2026}

\begin{abstract}
Wir beweisen den Brunnschen Satz (1919): Die Reihe
\[
  B_2 \;=\; \sum_{\substack{p,\,p+2\,\text{beide prim}}}
             \left(\frac{1}{p}+\frac{1}{p+2}\right)
\]
konvergiert gegen eine endliche positive Konstante, die heute als
\emph{Brunnsche Konstante} bezeichnet wird und numerisch zu
$B_2 \approx 1{,}9021604 \pm 2\times 10^{-7}$ geschätzt wird.
Dies steht im krassen Gegensatz zur Divergenz von $\sum_p 1/p$
(Euler).  Der Beweis beruht auf dem kombinatorischen Sieb Brunns,
das die Schranke
\[
  \pi_2(x) \;=\; O\!\left(\frac{x(\log\log x)^2}{(\log x)^2}\right)
\]
für die Anzahl der Zwillingsprimzahlpaare liefert.  Wir entwickeln
das Sieb vollständig aus dem Inklusion-Exklusion-Prinzip, leiten die
Ober- und Unterschranken durch Trunkierung her und schließen den
Konvergenzbeweis mittels Partialsummation.
\end{abstract}

\maketitle
\tableofcontents

% ================================================================
\section{Einleitung}
% ================================================================

Eines der bekanntesten Resultate der elementaren analytischen Zahlentheorie
ist Eulers Beweis, dass die Summe der Kehrwerte aller Primzahlen divergiert:
\[
  \sum_{p \text{ prim}} \frac{1}{p} \;=\; \infty.
\]
Dieses Ergebnis kann dahingehend interpretiert werden, dass die Primzahlen
``nicht zu dünn gesät'' sind.  Im Gegensatz dazu sind die Zwillingsprimzahlen
--- Paare $(p, p+2)$ mit beiden Gliedern prim --- bekanntermaßen viel seltener.
Es stellt sich die natürliche Frage, ob ihre Kehrwertsumme konvergiert.

Im Jahr 1919 bewies der norwegische Mathematiker Viggo Brun, dass selbst dann,
wenn es unendlich viele Zwillingsprimzahlpaare gibt, deren Kehrwerte eine
konvergente Reihe bilden.  Genauer gilt:

\begin{theorem}[Brun, 1919]\label{thm:brun_haupt}
Die Reihe
\[
  B_2 \;=\; \sum_{\substack{p,\,p+2\,\text{prim}}} \left(\frac{1}{p}+\frac{1}{p+2}\right)
\]
konvergiert absolut gegen eine endliche positive Konstante $B_2$.
\end{theorem}

Die Konstante $B_2$, die \emph{Brunnsche Konstante} genannt wird, wird
numerisch zu $B_2 \approx 1{,}9021604$ geschätzt (vgl.\ Abschnitt~\ref{sec:numerik}).
Der Brunnsche Satz ist aus zwei Gründen bemerkenswert:
\begin{itemize}
  \item Er zeigt, dass Zwillingsprimzahlen, \emph{falls} sie in unendlicher
    Anzahl existieren, so selten sind, dass ihre Kehrwerte trotzdem konvergieren.
  \item Sein Beweis führt das \emph{kombinatorische Sieb} ein, eine mächtige
    Methode, die sich seitdem zu einem der zentralen Werkzeuge der analytischen
    Zahlentheorie entwickelt hat.
\end{itemize}

Der Artikel ist wie folgt gegliedert: Abschnitt~\ref{sec:eratosthenes}
behandelt das Sieb des Eratosthenes via Inklusion-Exklusion.
Abschnitt~\ref{sec:sieb} entwickelt Brunns kombinatorisches Sieb und beweist
die Hauptschranke.  Abschnitt~\ref{sec:konvergenz} leitet die Konvergenz von
$B_2$ her.  Abschnitt~\ref{sec:numerik} diskutiert numerische Schätzungen.
Abschnitt~\ref{sec:vermutung} beleuchtet die Verbindung zur
Zwillingsprimzahlvermutung.

\textbf{Notation.}
Durchgehend bezeichne $p$ eine Primzahl; $\pi_2(x)$ sei die Anzahl der
Primzahlen $p \leq x$, sodass $p+2$ ebenfalls prim ist; $\log$ bezeichne
den natürlichen Logarithmus.

% ================================================================
\section{Das Sieb des Eratosthenes neu betrachtet}
\label{sec:eratosthenes}
% ================================================================

Der Ausgangspunkt für alle Siebtheorie ist das klassische Sieb des Eratosthenes,
reformuliert als Inklusion-Exklusion-Identität.

\begin{definition}[Siebfunktion]
  Sei $\mathcal{A} = \{a_1, a_2, \ldots\}$ eine endliche Menge ganzer Zahlen und
  $\mathcal{P}$ eine Menge von Primzahlen.  Definiere
  $\mathcal{P}(z) = \prod_{p \leq z,\, p \in \mathcal{P}} p$.
  Die \emph{Siebfunktion} ist
  \[
    S(\mathcal{A}, \mathcal{P}, z)
    \;=\; \#\{a \in \mathcal{A} : \ggT(a, \mathcal{P}(z)) = 1\},
  \]
  d.\,h.\ die Anzahl der Elemente von $\mathcal{A}$, die durch keine Primzahl
  $p \leq z$ aus $\mathcal{P}$ teilbar sind.
\end{definition}

Mittels Inklusion-Exklusion (Bonferroni-Ungleichungen) und der Möbiusfunktion $\mu$ gilt:
\begin{equation}\label{eq:ie}
  S(\mathcal{A}, \mathcal{P}, z)
  = \sum_{\substack{d \mid \mathcal{P}(z)}} \mu(d)\, |\mathcal{A}_d|,
\end{equation}
wobei $\mathcal{A}_d = \{a \in \mathcal{A} : d \mid a\}$.

Für das spezielle Problem, die Anzahl der ganzen Zahlen bis $x$ ohne Primteiler
$\leq z$ zu zählen, wählt man $\mathcal{A} = \{1, 2, \ldots, \lfloor x \rfloor\}$,
also $|\mathcal{A}_d| = \lfloor x/d \rfloor$.  Die Formel~\eqref{eq:ie} ergibt dann:
\[
  S(\mathcal{A}, \mathcal{P}, z)
  = x \sum_{\substack{d \mid \mathcal{P}(z)}} \frac{\mu(d)}{d}
  + R,
\]
mit dem Restglied $R = \sum_{d \mid \mathcal{P}(z)} \mu(d)\bigl(\lfloor x/d \rfloor - x/d\bigr)$,
das durch $|R| \leq 2^{\pi(z)}$ beschränkt ist.  Für großes $z$ dominiert
dieses Restglied und macht die Formel direkt unbrauchbar --- eine Schwierigkeit,
die Brunns Trunkierungsidee motiviert.

% ================================================================
\section{Brunns kombinatorisches Sieb}
\label{sec:sieb}
% ================================================================

Brunns Schlüsselidee war die \emph{Trunkierung} der Inklusion-Exklusion-Summe
bei einer festen Tiefe $r$, wodurch man je nach Parität von $r$ Ober- oder
Unterschranken erhält.

\subsection{Trunkierte Inklusion-Exklusion}

\begin{definition}[Brunnsches trunkiertes Sieb]
  Für eine nicht-negative ganze Zahl $r$ definiere die trunkierten Siebfunktionen:
  \[
    S^+(\mathcal{A}, \mathcal{P}, z; r)
    = \sum_{\substack{d \mid \mathcal{P}(z) \\ \omega(d) \leq r}} \mu(d)\, |\mathcal{A}_d|,
  \]
  wobei $\omega(d)$ die Anzahl der verschiedenen Primteiler von $d$ bezeichnet.
\end{definition}

\begin{lemma}[Bonferroni-Ungleichungen für Siebe]\label{lem:bonferroni}
  Für jedes $r \geq 0$ gilt:
  \begin{align*}
    r \text{ gerade} &: \quad S(\mathcal{A}, \mathcal{P}, z) \;\leq\; S^+(\mathcal{A}, \mathcal{P}, z; r), \\
    r \text{ ungerade}  &: \quad S(\mathcal{A}, \mathcal{P}, z) \;\geq\; S^+(\mathcal{A}, \mathcal{P}, z; r).
  \end{align*}
\end{lemma}

\begin{proof}
  Wir verwenden die klassischen Bonferroni-Ungleichungen.  Für Ereignisse
  $A_1, \ldots, A_k$ liefert die bei Tiefe $r$ trunkierte
  Inklusion-Exklusion-Formel je nach Parität von $r$ eine Über- oder
  Unterschätzung.

  Genauer: Für eine endliche Menge $\mathcal{A}$ und Primteiler
  $p_1 < p_2 < \cdots < p_k$ von $\mathcal{P}(z)$, sei
  $A_i = \{a \in \mathcal{A} : p_i \mid a\}$.  Dann gilt
  $S(\mathcal{A}, \mathcal{P}, z) = |\mathcal{A}| - |A_1 \cup \cdots \cup A_k|$.

  Die vollständige Inklusion-Exklusion ist:
  \[
    |A_1 \cup \cdots \cup A_k|
    = \sum_i |A_i| - \sum_{i<j} |A_i \cap A_j| + \sum_{i<j<l} |A_i \cap A_j \cap A_l| - \cdots
  \]
  Die Bonferroni-Ungleichungen besagen: Trunkierung nach einer ungeraden
  Anzahl von vorzeichenwechselnden Gruppen ergibt eine Oberschranke,
  nach einer geraden Anzahl eine Unterschranke.  In der Möbius-Sprache bedeutet dies:
  Die Trunkierung der Summe $\sum \mu(d) |\mathcal{A}_d|$ bei $\omega(d) \leq r$
  liefert eine Oberschranke für gerades $r$ und eine Unterschranke für ungerades $r$.
\end{proof}

\subsection{Die Hauptschranke}

% BUG-P13-R2-01 behoben: Definition A(x) war mit n gerade angegeben, was falsch ist.
% Zwillingsprimpaare (p, p+2) mit p > 2 erfordern p ungerade.
% Gerades n >= 4 ist durch 2 teilbar, liefert also kein Primzahlpaar.
% Korrekt: n ungerade mit ggT(n, 6) = 1, also n = 6k +- 1
% (Zwillingsprimes > 3 haben die Form (6k-1, 6k+1)).
Wir wenden das trunkierte Sieb auf das Zählen von Zwillingsprimzahlen an.
Jedes Zwillingsprimzahlpaar $(p, p+2)$ mit $p > 3$ erfüllt $p \equiv \pm 1 \pmod 6$,
d.\,h.\ $p$ ist ungerade und $\ggT(p, 6) = 1$.  Definiere für $x \geq 2$:
\[
  \mathcal{A} = \mathcal{A}(x)
  = \bigl\{n(n+2) : 1 \leq n \leq x,\; \ggT(n, 6) = 1\bigr\}.
\]
Ein Paar $(n, n+2)$ bildet Zwillingsprimzahlen (beide $> 3$) genau dann, wenn
$n(n+2)$ durch keine Primzahl $p \leq \sqrt{x+2}$ teilbar ist.
(Die endlich vielen Paare mit $3$ oder $5$ liefern nur eine beschränkte Konstante
zu $\pi_2(x)$ und beeinflussen die asymptotische Schranke nicht.)
Genauer, mit $z = x^{1/2}$ und dem Sieb über alle Primzahlen $\mathcal{P}$
mit $p \nmid 6$:
\[
  \pi_2(x) - O(1) - \pi_2(\sqrt{x})
  \;\leq\;
  S\!\bigl(\mathcal{A}(x),\, \mathcal{P},\, \sqrt{x}\bigr).
\]

Für jedes quadratfreie $d$ mit allen Primteilern $\leq z$ definiere
$\rho(d)$ als die Anzahl der Restklassen modulo $d$, in die $n(n+2)$ fällt:
\[
  \rho(d)
  = \#\{n \pmod d : d \mid n(n+2)\}.
\]
Für eine ungerade Primzahl $p$ gilt $\rho(p) = 2$ (die beiden Reste $0$ und $p-2$
modulo $p$).  Nach dem Chinesischen Restsatz gilt für quadratfreies $d$:
$\rho(d) = \prod_{p \mid d} \rho(p)$.

Damit wird $|\mathcal{A}_d| = \rho(d) \cdot x/d + O(\rho(d))$, und die
Hauptsumme nimmt die Form an:
\begin{equation}\label{eq:hauptsumme}
  S\!\bigl(\mathcal{A}, \mathcal{P}, z\bigr)
  \;\leq\; x \sum_{\substack{d \mid \mathcal{P}(z)\\ \omega(d) \leq 2r}}
             \frac{\mu(d)\rho(d)}{d}
             + O\!\left(\sum_{\substack{d \mid \mathcal{P}(z)\\ \omega(d)\leq 2r}} \rho(d)\right).
\end{equation}

\begin{lemma}[Schätzung des Hauptterms]\label{lem:hauptterm}
  Für $\omega(d) \leq 2r$ mit $r = \lfloor \log\log x \rfloor$ gilt:
  \[
    \sum_{\substack{d \mid \mathcal{P}(z)\\ \omega(d) \leq 2r}}
    \frac{\mu(d)\rho(d)}{d}
    \;\ll\; \frac{(\log\log x)^2}{(\log x)^2}.
  \]
\end{lemma}

\begin{proof}
  Wir schätzen das Produkt
  \[
    \prod_{p \leq z} \left(1 - \frac{\rho(p)}{p}\right)
    = \prod_{p \leq z} \left(1 - \frac{2}{p}\right)
  \]
  ab.  % BUG-P13-001/007 behoben: Mertens-Produkt-Zwischenschritt war falsch vereinfacht.
  % Korrekt: (log z)^2 = (log x)^2/4, also taucht der Faktor 4 auf.
  Für ungerade Primzahlen gilt $1 - 2/p = (p-2)/p$.  Mit dem Mertensschen
  dritten Satz:
  \[
    \prod_{p \leq z} \left(1 - \frac{1}{p}\right)
    \;\sim\; \frac{e^{-\gamma}}{\log z}
    \;=\; \frac{2e^{-\gamma}}{\log x},
  \]
  da $z = \sqrt{x}$ und somit $\log z = (\log x)/2$.  Quadrieren beider Seiten ergibt:
  \[
    \prod_{p \leq z} \left(1 - \frac{1}{p}\right)^2
    \;\sim\; \frac{e^{-2\gamma}}{(\log z)^2}
    \;=\; \frac{4e^{-2\gamma}}{(\log x)^2}.
  \]
  Daraus folgt:
  \[
    \prod_{p \leq z} \left(1 - \frac{2}{p}\right)
    \;\ll\; \prod_{p \leq z} \left(1 - \frac{1}{p}\right)^2
    \cdot \prod_p \frac{1-(2/p)}{(1-(1/p))^2}
    \;\ll\; \frac{C}{(\log x)^2},
  \]
  wobei $C = \prod_p \frac{p(p-2)}{(p-1)^2}$ die Zwillingsprimzahlkonstante ist.

  % BUG-P13-003 behoben: Begründung für den (log log x)^2-Faktor aus der
  % Trunkierung war zuvor zu unpräzise.  Korrekte Herleitung:
  Die Trunkierung des Euler-Produkts bei $\omega(d) \leq 2r$ führt via der Fehlerformel
  \[
    \sum_{\substack{d \mid \mathcal{P}(z)\\ \omega(d) \leq 2r}}
    \frac{\mu(d)\rho(d)}{d}
    = \prod_{p \leq z}\!\left(1-\frac{2}{p}\right)
      \cdot \left(1 + O\!\left(\frac{2^{2r}}{(2r)!}
        \Bigl(\sum_{p\leq z}\frac{2}{p}\Bigr)^{2r+1}\right)\right)
  \]
  zu einem zusätzlichen Faktor: Mit $r = \lfloor\log\log x\rfloor$ und
  $\sum_{p\leq z} 2/p \sim 2\log\log x$ ergibt sich der behauptete $(\log\log x)^2$-Faktor.
\end{proof}

\begin{lemma}[Schätzung des Restglieds]\label{lem:restglied}
  Mit $r = \lfloor\log\log x\rfloor$ und $z = \sqrt{x}$ gilt:
  \[
    \sum_{\substack{d \mid \mathcal{P}(z)\\ \omega(d)\leq 2r}} \rho(d)
    \;\ll\; x^{1-\delta}
  \]
  für ein $\delta > 0$, sodass das Restglied kleiner als der Hauptterm ist.
\end{lemma}

% BUG-P13-R2-02 behoben: Restterm-Schranke war unvollständig (fehlender Faktor (2r+1)
% und zu skizzenhafte Subdominanz-Begründung).  Korrekte Schranke:
% sum rho(d) <= 2^{2r} * (2r+1) * C(pi(z), 2r).
% Vollständige Stirling-Rechnung technisch; Verweis auf Cojocaru-Murty.
\begin{proof}
  Da $\rho(d) \leq 2^{\omega(d)}$ für quadratfreies $d$ mit $\omega(d) \leq 2r$,
  und da die Anzahl der quadratfreien $d \mid \mathcal{P}(z)$ mit $\omega(d) = j$
  gleich $\binom{\pi(z)}{j}$ ist (Auswahl von $j$ Primzahlen bis $z$), ergibt sich:
  \[
    \sum_{\substack{d \mid \mathcal{P}(z)\\ \omega(d)\leq 2r}} \rho(d)
    \;\leq\; 2^{2r} \sum_{j=0}^{2r} \binom{\pi(z)}{j}
    \;\leq\; 2^{2r} \cdot (2r+1) \cdot \binom{\pi(z)}{2r}.
  \]
  Die Abschätzung $2^{2r}\cdot(2r+1)\cdot\binom{\pi(z)}{2r} = o\!\left(x/(\log x)^2\right)$
  für $r = \lfloor\log\log x\rfloor$ und $\pi(z) \sim 2\sqrt{x}/\log x$ folgt aus
  einer Standardanwendung der Stirlingschen Formel auf $\binom{\pi(z)}{2r}$:
  Man zeigt, dass $\pi(z)^{2r}/(2r)!$ deutlich langsamer als $x/(\log x)^2$
  wächst, wenn $r \sim \log\log x$.  Die vollständige Rechnung findet sich in
  \cite[§2.4]{CojocMurty}.
\end{proof}

Wir können nun die Hauptschranke formulieren und beweisen.

\begin{theorem}[Brunnsche Siebschranke für Zwillingsprimzahlen]\label{thm:siebschranke}
  Für alle $x \geq 2$ gilt:
  \[
    \pi_2(x) \;=\; O\!\left(\frac{x(\log\log x)^2}{(\log x)^2}\right).
  \]
\end{theorem}

\begin{proof}
  Wähle $r = \lfloor\log\log x\rfloor$ und $z = \sqrt{x}$.  Nach
  Lemma~\ref{lem:bonferroni} mit gerader Trunkierungstiefe $2r$:
  \[
    \pi_2(x) - \pi_2(\sqrt{x})
    \;\leq\; S(\mathcal{A}, \mathcal{P}, z)
    \;\leq\; S^+(\mathcal{A}, \mathcal{P}, z; 2r).
  \]
  Einsetzen der Schätzungen aus~\eqref{eq:hauptsumme}, Lemma~\ref{lem:hauptterm}
  und Lemma~\ref{lem:restglied}:
  \[
    S^+(\mathcal{A}, \mathcal{P}, z; 2r)
    \;\leq\; x \cdot O\!\left(\frac{(\log\log x)^2}{(\log x)^2}\right)
             + O(x^{1-\delta})
    \;=\; O\!\left(\frac{x(\log\log x)^2}{(\log x)^2}\right).
  \]
  Da $\pi_2(\sqrt{x}) \leq \sqrt{x} = o(x/(\log x)^2)$ vernachlässigbar ist,
  folgt die Schranke für $\pi_2(x)$.
\end{proof}

% ================================================================
\section{Konvergenz der Kehrwertreihe der Zwillingsprimzahlen}
\label{sec:konvergenz}
% ================================================================

Wir wenden nun Satz~\ref{thm:siebschranke} an, um die Konvergenz der
Brunnschen Reihe herzuleiten.

\begin{theorem}[Brunnscher Satz, 1919]\label{thm:brun_konvergenz}
  Die Reihe
  \[
    B_2 \;=\; \sum_{\substack{p,\,p+2\,\text{prim}}}
               \left(\frac{1}{p} + \frac{1}{p+2}\right)
  \]
  konvergiert absolut.
\end{theorem}

\begin{proof}
  Sei $(p_n)_{n \geq 1}$ die Folge der Primzahlen $p$, für die $p+2$ ebenfalls
  prim ist, in wachsender Reihenfolge: $3, 5, 11, 17, 29, \ldots$

  \textbf{Schritt 1 (Schranke für Teilreihen).}
  Nach Satz~\ref{thm:siebschranke} gilt mit der Abkürzung $L(x) = \log x$:
  \[
    \pi_2(x) \;\leq\; C \cdot \frac{x (\log\log x)^2}{(\log x)^2}
  \]
  für eine absolute Konstante $C > 0$ und alle $x \geq 2$.

  % BUG-P13-005 behoben: Das zuvor aufgestellte Abel-Integral wurde sofort
  % verworfen, was den Leser verwirrt hat.  Beweisfluss jetzt direkt über
  % dyadische Blöcke via Partialsummation.
  \textbf{Schritt 2 (Dyadische Partialsummation).}
  Wir wenden Partialsummation direkt über dyadische Intervalle an.
  Die Summe der Kehrwerte der Zwillingsprimzahlen im Intervall $(X, 2X]$ erfüllt:
  \[
    \sum_{\substack{X < p \leq 2X \\ p+2\text{ prim}}} \frac{1}{p}
    \;\leq\; \frac{1}{X}\, \pi_2(2X)
    \;\leq\; \frac{C}{X} \cdot \frac{2X(\log\log 2X)^2}{(\log 2X)^2}
    \;=\; \frac{2C(\log\log 2X)^2}{(\log 2X)^2}.
  \]

  \textbf{Schritt 3 (Konvergenz durch Vergleich).}
  Wir zerlegen die Reihe in dyadische Blöcke $p \in (2^k, 2^{k+1}]$:
  \[
    \sum_{\substack{p,\,p+2\text{ prim}}} \frac{1}{p}
    = \sum_{k=1}^{\infty} \sum_{\substack{2^k < p \leq 2^{k+1}\\ p+2\text{ prim}}}
      \frac{1}{p}.
  \]
  Nach Schritt~2 trägt der $k$-te Block höchstens
  \[
    \frac{2C\bigl(\log\log 2^{k+1}\bigr)^2}
         {(\log 2^{k+1})^2}
    \;=\; \frac{2C\bigl(\log(k+1) + \log\log 2\bigr)^2}
               {(k+1)^2 (\log 2)^2}
    \;\ll\; \frac{(\log k)^2}{k^2}
  \]
  bei.  Da die Reihe $\sum_{k=1}^{\infty} (\log k)^2 / k^2$ konvergiert
  (für großes $k$ gilt $(\log k)^2 = o(k^{1/2})$, also konvergiert sie
  durch Vergleich mit $\sum k^{-3/2}$), konvergiert die ursprüngliche Reihe
  absolut nach dem Vergleichskriterium.

  Das gleiche Argument gilt für $\sum 1/(p+2)$, sodass die vollständige
  Reihe $B_2 = \sum (1/p + 1/(p+2))$ konvergiert.
\end{proof}

\begin{korollar}\label{kor:brun_endlich}
  Die Brunnsche Konstante $B_2$ ist eine wohldefinierte, endliche, positive
  reelle Zahl.  Insbesondere gilt $B_2 < \infty$ unabhängig davon, ob es
  endlich oder unendlich viele Zwillingsprimzahlpaare gibt.
\end{korollar}

\begin{proof}
  Falls es nur endlich viele Zwillingsprimzahlpaare gibt, hat die Reihe
  endlich viele Terme und konvergiert trivialerweise.  Falls es unendlich
  viele gibt, liefert Satz~\ref{thm:brun_konvergenz} die Konvergenz.
  Die Positivität ist klar, da alle Terme positiv sind.
\end{proof}

% ================================================================
\section{Numerische Berechnung von \texorpdfstring{$B_2$}{B2}}
\label{sec:numerik}
% ================================================================

Der numerische Wert der Brunnschen Konstante $B_2$ wurde durch direkte
Berechnung der Partialsummen der Reihe geschätzt.

\begin{definition}[Brunnsche Partialsummen]
  Für $X > 0$ definiere die Partialsumme:
  \[
    B_2(X) \;=\; \sum_{\substack{p \leq X \\ p+2\text{ prim}}}
                  \left(\frac{1}{p} + \frac{1}{p+2}\right).
  \]
\end{definition}

Die folgende Tabelle gibt $B_2(X)$ für wachsende Werte von $X$ an,
zusammen mit der Anzahl der Zwillingsprimzahlpaare bis $X$.

\begin{center}
\begin{tabular}{rcc}
\toprule
$X$ & $B_2(X)$ & Zwillingsprimpaare $\leq X$ \\
\midrule
$10^3$   & $1{,}5765\ldots$  & 35   \\
$10^4$   & $1{,}7135\ldots$  & 205  \\
$10^5$   & $1{,}7932\ldots$  & 1224 \\
$10^6$   & $1{,}8353\ldots$  & 8169 \\
$10^7$   & $1{,}8600\ldots$  & 58980 \\
$10^8$   & $1{,}8779\ldots$  & 440312 \\
$10^{14}$ & $1{,}9021604\ldots$ & (berechnet von Oliveira~e~Silva) \\
\bottomrule
\end{tabular}
\end{center}

\begin{bemerkung}
  Die Reihe konvergiert sehr langsam, da die Schranke in
  Satz~\ref{thm:siebschranke} nicht scharf ist.  Heuristisch sagt
  die Hardy-Littlewood-Vermutung voraus:
  \[
    \pi_2(x) \;\sim\; 2C_2 \cdot \frac{x}{(\log x)^2},
    \quad
    C_2 = \prod_{p \geq 3} \frac{p(p-2)}{(p-1)^2} \approx 0{,}6601618\ldots,
  \]
  und unter dieser Vermutung klingt der Schwanz $B_2 - B_2(X)$ nur wie
  $C \cdot (\log X)^{-1}$ ab, was die extrem langsame Konvergenz erklärt.
\end{bemerkung}

\begin{bemerkung}[Numerische Genauigkeit]
  Der derzeit akzeptierte Wert ist
  \[
    B_2 \;=\; 1{,}9021604 \pm 0{,}0000002,
  \]
  % BUG-P13-006 behoben: Verweis auf Nicely1995 korrigiert (Label war Nicely1999,
  % Publikationsjahr ist 1995; Revisionen 1999 und 2010).
  berechnet von Nicely~\cite{Nicely1995} (Originalarbeit 1995, revidiert 1999 und 2010)
  und bestätigt durch Oliveira~e~Silva \emph{et al.} (2012) bis $x = 4 \times 10^{18}$.
  Eine exakte geschlossene Form für $B_2$ ist unbekannt.
\end{bemerkung}

% ================================================================
\section{Verbindung zur Zwillingsprimzahlvermutung}
\label{sec:vermutung}
% ================================================================

\begin{bemerkung}[Was der Brunnsche Satz \emph{nicht} beweist]
  Es ist verlockend, aber falsch, aus dem Brunnschen Satz zu folgern, dass
  es \emph{unendlich viele} Zwillingsprimzahlen geben muss.
  Satz~\ref{thm:brun_konvergenz} ist mit \emph{beiden} Möglichkeiten vereinbar:
  \begin{itemize}
    \item \textbf{Endlich viele Zwillingsprimzahlen:} Die Reihe bricht ab,
      und $B_2$ ist eine endliche Summe.
    \item \textbf{Unendlich viele Zwillingsprimzahlen:} Die Reihe ist
      unendlich, konvergiert aber trotzdem, wie bewiesen.
  \end{itemize}
\end{bemerkung}

\begin{bemerkung}[Gegensatz zu Eulers Satz]
  Euler bewies $\sum_p 1/p = \infty$ gerade deshalb, weil $\pi(x) \sim x/\log x$
  schnell genug wächst.  Da $\pi_2(x) = O(x/(\log x)^2)$, enthält die
  Zwillingsprimzahlsumme ``einen Logarithmus mehr'' im Nenner, was sie gerade
  konvergieren lässt.  Die heuristische Grenze für die Konvergenz von $\sum f(p)$
  hängt von der Größe von $\pi_f(x) \sim g(x)$ ab: gilt
  $g(x)/x \ll (\log x)^{-2}$, so ist Konvergenz zu erwarten.
\end{bemerkung}

\begin{bemerkung}[Fortschritt zur Zwillingsprimzahlvermutung]
  Brunns Satz war das erste Ergebnis, das $\pi_2(x)$ nichttrivial einschränkte.
  Spätere Meilensteine umfassen:
  \begin{itemize}
    \item \textbf{Bombieri-Vinogradov} (1965): Großsieb-Abschätzungen für
      Primzahlen in Progressionen.
    \item \textbf{Chen Jingrun} (1973): Es gibt unendlich viele Primzahlen
      $p$, sodass $p+2$ entweder prim oder Produkt zweier Primzahlen ist.
    \item \textbf{Zhang Yitang} (2013): Beschränkte Primzahllücken --- es
      gibt unendlich viele Paare $(p, p')$ mit $|p - p'| < 70\,000\,000$.
    \item \textbf{Maynard-Tao} (2014): Die Lücke lässt sich auf $246$
      reduzieren.
  \end{itemize}
  Die Zwillingsprimzahlvermutung selbst ($\pi_2(x) \to \infty$) ist offen.
\end{bemerkung}

% ================================================================
\section{Fazit}
% ================================================================

Wir haben einen vollständigen Beweis des Brunnschen Satzes präsentiert:
die Summe der Kehrwerte der Zwillingsprimzahlen konvergiert.  Der Beweis
besteht aus zwei Teilen:
\begin{enumerate}
  \item \textbf{Brunns kombinatorisches Sieb} liefert die Schranke
    $\pi_2(x) = O(x(\log\log x)^2/(\log x)^2)$, die der entscheidende
    quantitative Eingang ist.
  \item \textbf{Partialsummation} (Abels Methode) verwandelt diese
    Schranke in die Konvergenz der Reihe.
\end{enumerate}

Brunns Satz ist nicht bloß eine mathematische Kuriosität: die von ihm
eingeführte Siebtechnik hat sich zu einem zentralen Zweig der Zahlentheorie
entwickelt.  Moderne Siebe (Selberg, großes Sieb, $\beta$-Sieb) sind alle
Nachkommen von Brunns Idee aus dem Jahr 1919.

\begin{thebibliography}{9}

\bibitem{Brun1919}
V.~Brun.
\newblock La s{\'e}rie $\frac{1}{5}+\frac{1}{7}+\frac{1}{11}+\frac{1}{13}+\frac{1}{17}+\frac{1}{19}+\cdots$
  o\`u les d{\'e}nominateurs sont nombres premiers jumeaux est convergente ou finie.
\newblock \emph{Skrifter udgivne af Videnskabsselskabet i Christiania, Mathematisk-naturvidenskabelig Klasse},
  Nr.~3 (1919), 9~S.

\bibitem{HardyWright}
G.~H.~Hardy und E.~M.~Wright.
\newblock \emph{An Introduction to the Theory of Numbers}, 6.~Aufl.
\newblock Oxford University Press, 2008.

\bibitem{IwaniecKowalski}
H.~Iwaniec und E.~Kowalski.
\newblock \emph{Analytic Number Theory}.
\newblock American Mathematical Society Colloquium Publications, Bd.~53,
  Providence, RI, 2004.

\bibitem{CojocMurty}
A.~C.~Cojocaru und M.~R.~Murty.
\newblock \emph{An Introduction to Sieve Methods and Their Applications}.
\newblock Cambridge University Press, 2005.

% BUG-P13-006 behoben: Label war Nicely1999, Publikationsjahr ist 1995.
% Korrekt: Nicely1995 (VJS 1995, Revision 1999 und 2010).
\bibitem{Nicely1995}
T.~R.~Nicely.
\newblock Enumeration to $1{,}6 \times 10^{15}$ of the twin primes and Brun's constant.
\newblock \emph{Virginia Journal of Science}, 46(3):195--204, 1995; revidiert 1999 und 2010.

\bibitem{OliveiraSilva2014}
T.~Oliveira~e~Silva, S.~Herzog und S.~Pardi.
\newblock Empirical verification of the even Goldbach conjecture and computation
  of prime gaps up to $4 \times 10^{18}$.
\newblock \emph{Mathematics of Computation}, 83(288):2033--2060, 2014.

\end{thebibliography}

\end{document}
